\section{Introduction}

Las heladas meteorológicas, definidas como el descenso de la
temperatura del aire por debajo de $0^{\circ}\text{C}$, constituyen uno de los
fenómenos climáticos extremos de mayor recurrencia en las zonas altoandinas del
Perú. En un país donde amplios sectores de la población rural dependen de una
agricultura y ganadería de subsistencia altamente sensibles a la variabilidad
térmica nocturna, estos eventos representan una amenaza estructural para la
seguridad alimentaria y la estabilidad económica de las comunidades más
vulnerables. La creciente irregularidad de los patrones climáticos ha
intensificado la frecuencia y severidad de estos episodios, situando su
predicción anticipada como un problema de relevancia nacional.

Desde una perspectiva socioeconómica, la ocurrencia de una helada puede destruir
en cuestión de horas la totalidad de un cultivo sensible, provocando pérdidas
directas para el productor y presionando las cadenas de abastecimiento
regionales. La capacidad de anticipar estos eventos con antelación suficiente
permite activar medidas de mitigación en campo, reduciendo el impacto económico
y protegiendo los medios de vida de las poblaciones afectadas. Técnicamente, sin
embargo, los modelos numéricos de predicción meteorológica tradicionales suelen
carecer de la resolución espacial necesaria para capturar los microclimas
locales que determinan la ocurrencia de heladas a escala puntual. En este
contexto, los métodos de aprendizaje automático supervisado, entrenados
directamente sobre registros históricos de estaciones meteorológicas, han
emergido como una alternativa computacionalmente eficiente y de bajo costo para
la predicción de estos eventos a nivel local
\cite{diedrichs2018,talsma2023,rozante2023,perez2025}.

No obstante, la aplicación efectiva de estas técnicas sobre datos meteorológicos
reales impone desafíos que trascienden la simple selección de un algoritmo. En
términos de minería de datos, la predicción de heladas se formula como un
problema de \emph{clasificación binaria supervisada fuertemente desbalanceado},
en el que la clase de interés (ocurrencia de helada) es estadísticamente
minoritaria frente a las condiciones climáticas normales. Este desbalance sesga
a los clasificadores convencionales hacia la clase mayoritaria, tornando
inadecuadas métricas como la exactitud y exigiendo estrategias explícitas de
mitigación \cite{azhar2023,wang2021}. A este desafío se suman tres
complicaciones inherentes al dominio: (i) la presencia de códigos de error de
instrumentación no documentados que contaminan la topología estadística de los
datos crudos; (ii) el riesgo de fuga de información (\emph{data leakage}) al
derivar la variable objetivo a partir de una medición física que, de conservarse
entre los predictores, induciría una regla determinista trivial; y (iii) la alta
multicolinealidad entre variables climáticas interdependientes, que introduce
redundancia y degrada la convergencia del modelado. El presente trabajo aborda
integralmente estos cuatro aspectos mediante un pipeline reproducible diseñado
específicamente para el contexto meteorológico peruano.

Las principales contribuciones de este trabajo son:
\begin{enumerate}
\item Se identifica y corrige de forma sistemática un código de error de sensor
(\texttt{-999}) no documentado en el conjunto de datos crudo, mediante un
diagnóstico exploratorio comparativo de las funciones de densidad antes y después
de la corrección, recuperando la topología estadística real de las variables.
\item Se diseña una variable objetivo físicamente fundamentada
(\emph{Target\_Helada}) a partir del umbral termodinámico de $0^{\circ}\text{C}$,
eliminando explícitamente el predictor definitorio del espacio de características
para prevenir la fuga de información por predicción tautológica.
\item Se construye un pipeline de preprocesamiento robusto a valores atípicos
físicamente válidos, combinando imputación por mediana y \texttt{RobustScaler},
justificado matemáticamente frente a la estandarización clásica y a la
eliminación de \emph{outliers}, seguido de una reducción de dimensionalidad
mediante Análisis de Componentes Principales (PCA) que retiene el 95\% de la
varianza empírica.
\item Se optimiza un \emph{Random Forest Classifier} con ponderación balanceada
de clases mediante \texttt{GridSearchCV}, empleando el F1-Score Macro como función
objetivo para penalizar directamente los falsos negativos críticos, en lugar de
recurrir al sobremuestreo sintético empleado por trabajos previos.
\item Se realiza una evaluación integral del modelo mediante métricas robustas al
desbalance de clases (curva ROC-AUC, curva Precision-Recall e índice Kappa de
Cohen), acompañada de un análisis de la capacidad discriminante de los
componentes principales.
\end{enumerate}

El resto de este artículo está organizado de la siguiente manera. La
Sección~\ref{sec:relwork} presenta los trabajos relacionados y una comparación
sistemática con el enfoque propuesto. La Sección~\ref{sec:dataset} describe el
conjunto de datos y el análisis exploratorio, incluyendo la detección de
anomalías. La Sección~\ref{sec:preprocessing} detalla el preprocesamiento
matemático de los datos. La Sección~\ref{sec:methodology} expone la metodología
propuesta, sus algoritmos y la estrategia de optimización. La
Sección~\ref{sec:implementation} describe la implementación de software y las
consideraciones de reproducibilidad. La Sección~\ref{sec:results} presenta los
resultados experimentales. La Sección~\ref{sec:discussion} interpreta los hallazgos, sus
fortalezas y limitaciones. Finalmente, la Sección~\ref{sec:conclusions} presenta
las conclusiones y el trabajo futuro.